%%%%%%%%%%%%%%%%%%%%%%%%%%%%%%%%%%%%%%%%%%%%%%%%%%%%%%%%%%
%%
%%  PROJECT: prs900template
%%
%%  Created by Etienne van Delden on 24-05-11.
%%  Copyright 2011 __MyCompanyName__. 
%%  
%%  All rights reserved.
%%
%%%%%%%%%%%%%%%%%%%%%%%%%%%%%%%%%%%%%%%%%%%%%%%%%%%%%%%%%%

\documentclass[9pt]{memoir}

%%%%%%%%%%%%%%%%%%%%%%%%%%%%%%%%%%%%%%%%%%%%%%%%%%%%%%%%%%
%% 
%% MARK: Packages declarations
%%
%%%%%%%%%%%%%%%%%%%%%%%%%%%%%%%%%%%%%%%%%%%%%%%%%%%%%%%%%%



\usepackage{amssymb}
\usepackage{amsmath}
\usepackage[english]{babel}
\usepackage{lipsum}

\usepackage{XuBu}

%%%%%%%%%%%%%%%%%%%%%%%%%%%%%%%%%%%%%%%%%%%%%%%%%%%%%%%%%%
%% 
%% MARK: Input macro file
%% 
%%%%%%%%%%%%%%%%%%%%%%%%%%%%%%%%%%%%%%%%%%%%%%%%%%%%%%%%%%

\renewcommand{\cover}{\includegraphics[width=5cm]{cover}}


%%%%%%%%%%%%%%%%%%%%%%%%%%%%%%%%%%%%%%%%%%%%%%%%%%%%%%%%%%
%% 
%% Begin Document
%% 
%%%%%%%%%%%%%%%%%%%%%%%%%%%%%%%%%%%%%%%%%%%%%%%%%%%%%%%%%%

\begin{document}
	
	%%%%%%%%%%%%%%%%%%%%%%%%%%%%%%%%%%%%%%%%%%%%%%%%%%%%%%%%%%
	%% 
	%% MARK: Title & Data
	%% 
	%%%%%%%%%%%%%%%%%%%%%%%%%%%%%%%%%%%%%%%%%%%%%%%%%%%%%%%%%%
    
    \renewcommand{\contentsname}{Table of Contents} %Rename ToC to Table of Contents
    
	\title{Extra long title for prs900template template}
    \subtitle{Using extra \LaTeX awesomeness!}
	\author{Etienne van Delden}
	\date{\today}
    %\titlepic{\includegraphics[height=8.2cm]{cover.png}}
    
    
   	\maketitle
    \thispagestyle{empty}
    
	%%%%%%%%%%%%%%%%%%%%%%%%%%%%%%%%%%%%%%%%%%%%%%%%%%%%%%%%%%
	%% 
	%% MARK: Abstract
	%% 
	%%%%%%%%%%%%%%%%%%%%%%%%%%%%%%%%%%%%%%%%%%%%%%%%%%%%%%%%%%
    \begin{figure}[b]
        \begin{abstract}
            The Abstract
        \end{abstract}
    \end{figure}

    \newpage
   \setcounter{tocdepth}{1}
    \tableofcontents*
    
    
    %\pagestyle{mystyle}
        
	%%%%%%%%%%%%%%%%%%%%%%%%%%%%%%%%%%%%%%%%%%%%%%%%%%%%%%%%%%
	%% 
	%% MARK: Introduction
	%% 
	%%%%%%%%%%%%%%%%%%%%%%%%%%%%%%%%%%%%%%%%%%%%%%%%%%%%%%%%%%
	\mainmatter
	\chapter{Introduction}
        \chapterprecis{The chapterprecis, an awesome standard \LaTeX command}
	
        \section{The new section}\noindent    
        A read aloud text command!
        \readaloud{When the players arive, read:}{\lipsum[1-1]}
         \lipsum[2-2]
        \section{the next section}
            \lipsum[3-4]
        \subsection{last section}
            \lipsum[4-5]
	
	%%%%%%%%%%%%%%%%%%%%%%%%%%%%%%%%%%%%%%%%%%%%%%%%%%%%%%%%%%
	%% 
	%% MARK: The chapters
	%% 
	%%%%%%%%%%%%%%%%%%%%%%%%%%%%%%%%%%%%%%%%%%%%%%%%%%%%%%%%%%
	
	\chapter{The longer chapter title that should fill at least several lines of text} 
	\label{cha:}
    \chapterprecis{A slightly longer chapterprecis, but only because this sentence contains words that have no use what so over. Filler words to not fully rely on Lorem Ipsum}
        \lipsum[1-6]

	
	%%%%%%%%%%%%%%%%%%%%%%%%%%%%%%%%%%%%%%%%%%%%%%%%%%%%%%%%%%
	%% 
	%% MARK: Bibliography
	%% 
	%%%%%%%%%%%%%%%%%%%%%%%%%%%%%%%%%%%%%%%%%%%%%%%%%%%%%%%%%%
	
	\begin{thebibliography}{XXXX}
		
		\bibitem[XXXX]{XXXX}
		\newblock{Name,}
		\newblock{\em Title,}
		\newblock{ vol. X.}
		\newblock{ pp. XX--YY.}
		\newblock{Editor, Town,}
		\newblock{Year.}
		
	\end{thebibliography}
	
\end{document}